Esta guía resume un recorrido recomendado para presentar la aplicación ante un tribunal. No sustituye al manual de usuario, sino que propone una secuencia breve y ordenada para enseñar las funcionalidades principales sin perder la continuidad narrativa del proyecto.

\section{Preparación previa}

Antes de iniciar la demostración conviene comprobar:

\begin{itemize}
    \item que la aplicación está arrancada en el puerto previsto y que la base de datos contiene datos de prueba suficientes;
    \item que existen credenciales disponibles para los perfiles \texttt{admin}, \texttt{commercial} y \texttt{client};
    \item que el navegador tiene permisos adecuados para cámara o ubicación si se desea enseñar lectura QR o geolocalización;
    \item que el PDF de la memoria y las capturas principales están accesibles por si se necesita apoyar alguna explicación técnica.
\end{itemize}

\section{Recorrido recomendado}

\begin{enumerate}
    \item \textbf{Inicio de sesión y roles.} Abrir \path{/login}, iniciar sesión y explicar que el sistema redirige al panel correspondiente según el rol. Si hay tiempo, mostrar brevemente que las rutas protegidas no son accesibles sin sesión.
    \item \textbf{Panel administrador.} Desde el panel de administración, enseñar la gestión de usuarios, solicitudes de registro y clientes. Conviene remarcar que el administrador centraliza la operación y que las acciones sensibles quedan auditadas.
    \item \textbf{Auditoría y ajustes.} Abrir \path{/admin/audit} y \path{/admin/settings}. Mostrar la trazabilidad administrativa, la exportación \texttt{CSV} y la configuración visible de canales de avisos. Conviene explicar que el rate limiting, los inventarios técnicos y la base de Factusol/n8n quedan implementados como infraestructura interna, pero no como pantallas de uso diario.
    \item \textbf{Clientes y asignaciones comerciales.} Mostrar un cliente profesional, su información de contacto, ubicación y asignación a comercial. Este punto conecta el módulo de usuarios con la operativa diaria.
    \item \textbf{Rutas y visitas.} Entrar con un comercial, abrir \path{/commercials/routes} y una visita concreta. Explicar la planificación diaria, la geocodificación, las estimaciones y la integración del reparto dentro de visitas comerciales.
    \item \textbf{Catálogo y coloración.} Mostrar el catálogo de productos, filtros, ficha de producto, recursos asociados y cartas de color. Si se enseña un tinte, destacar la relación entre producto pedible y tonalidad cromática.
    \item \textbf{Pedido, reparto y cobro mínimo.} Crear o abrir un pedido, revisar sus líneas, confirmar el flujo de entrega con QR si hay datos preparados y enseñar el registro de cobro mínimo cuando el pedido está entregado.
    \item \textbf{Área de cliente.} Acceder como cliente profesional para mostrar catálogo, pedidos propios, información de reparto y perfil. Este punto evidencia que el sistema no es solo administrativo.
    \item \textbf{Gestión técnica del salón.} Abrir \path{/clients/salon-clients}, mostrar una ficha de cliente del salón, registrar o revisar un servicio, tonalidades, imágenes finales, plantillas y borrador técnico de correo.
    \item \textbf{Comunicaciones multicanal.} Mostrar promociones, formaciones, avisos, activación Push PWA y descuentos aplicados a pedido. Conviene explicar que el canal interno ya es la fuente principal y que correo SMTP y Web Push están implementados como canales opcionales, mientras que \textcolor{red}{\textit{WhatsApp} y los envíos planificados en segundo plano quedan como trabajo futuro}.
\end{enumerate}

\section{Cierre de la demo}

Para cerrar la presentación puede resumirse el recorrido en tres ideas:

\begin{itemize}
    \item la aplicación cubre un flujo operativo completo desde administración hasta cliente final profesional;
    \item los módulos se han construido de forma incremental, con persistencia, roles, trazabilidad y validaciones reales;
    \item \textcolor{red}{las líneas futuras no corrigen carencias básicas, sino que amplían una base ya funcional hacia automatización, integraciones externas y explotación avanzada de datos.}
\end{itemize}
